\documentclass[11pt]{article}
\usepackage[margin=25mm]{geometry}
\usepackage{fontspec}
\setmainfont{Latin Modern Roman}
\usepackage{amsmath,amssymb,mathtools}
\usepackage{unicode-math}
\setmathfont{Latin Modern Math}
\usepackage{xurl,hyperref}
\hypersetup{colorlinks=true,urlcolor=blue}
\setcounter{secnumdepth}{0}
\setlength{\parindent}{0pt}
\setlength{\parskip}{5pt}
\setlength{\emergencystretch}{3em}
\newcommand{\sref}[2]{\hyperlink{source:#1:#2}{[S#2]}}
\newcommand{\cataloguescope}[1]{\paragraph{Recorded target.}#1}
\newcommand{\R}{\mathbb R}
\newcommand{\C}{\mathbb C}
\newcommand{\N}{\mathbb N}
\newcommand{\Z}{\mathbb Z}
\newcommand{\Q}{\mathbb Q}
\newcommand{\F}{\mathbb F}
\newcommand{\eps}{\varepsilon}
\begin{document}
% UnsolvedMath ID 8100001; source catalogue code AMR-080-0001
\setcounter{section}{8100000}
\section{KdV with almost periodic initial data}\label{8100001}
{\small UnsolvedMath ID 8100001 · Mathematical Physics}
\subsection{Definitions and mathematical statement}

\paragraph{Shared notation.}$\mathbb N=\{1,2,\ldots\}$; the standard number fields have their usual meanings. All additional conventions are stated locally.


\paragraph{Definitions and notation.}A bounded function is Bohr almost periodic if its translates form a precompact family in the uniform norm; equivalently it is uniformly approximable by trigonometric polynomials. Quasiperiodic data depend on finitely many torus frequencies. The KdV normalization used here is $u_t+uu_x+u_{xxx}=0$. Almost periodic evolution in time requires a specified function-space topology, not merely boundedness of a solution.

\paragraph{Catalogue wording retained for traceability.}
Consider the Korteweg--de Vries (KdV) equation \[u_t+uu_x+u_{xxx}=0,\qquad u(x,0)=u_0(x),\qquad x\in\mathbb R.\] In the 1970's, McKean and Trubowitz proved the remarkable result that if the initial data $u_0$ is periodic, $u_0(x+p)= u_0(x)$ for some $p>0$, then the solution $u(x,t)$ of this equation is almost periodic in time. This result leads to the following natural conjecture: The same is true if $u_0(x)$ is almost periodic, i.e., if the initial data is almost periodic in space, the solution evolves almost periodically in time. This is a very challenging problem with many new and unique features. For example, for $u_0(x)$ almost periodic, one does not know, by standard PDE methods, whether a solution $u(x,t)$ exists, even for a small time, never mind globally. In spectacular and seminal work over the last few years, D. Damanik and M. Goldstein, joined later by I. Binder and M. Lukic, have partially resolved this conjecture in the affirmative, namely, for small quasi-periodic analytic initial data with Diophantine frequencies, unique global solutions exist and are almost periodic in time. The key insight of Damanik et al.\ was how to utilize the formal integrability of KdV in an effective analytical manner. Not only are the arguments in Damanik et al.\ applicable to other (formally) integrable systems (for example proves a Toda lattice analog of the KdV theorem), but some of their arguments are also applicable to non-integrable systems. It is a natural, open and very challenging conjecture that the result of Damanik et al.\ remains true for (suitable) perturbations of KdV, i.e., for (suitable) perturbations of KdV, the Cauchy problem with (suitable) almost periodic initial data $u_0(x)$, has a unique global solution $u(x,t)$, which evolves almost periodically in time. After all, the finite-dimensional analog of this result is precisely the content of classical KAM theory. We note that such results for perturbations of KdV and the nonlinear Schr\"{o}dinger equation (NLS) with periodic initial data $u_0(x)$, have already been proved by S. Kuksin, T. Kappeler and others. With the discovery of quasi-crystals, one can anticipate, in particular, that interest in PDE problems on quasi-periodic backgrounds, will surely grow.
\paragraph{Scope and source qualification.}The original source records partial progress for small analytic quasiperiodic data with Diophantine frequencies. Its words suitable perturbations and suitable data are part of a research programme; no unsupported universal perturbation class is imposed.
\subsection{Short English statement}
Prove global and time-almost-periodic evolution for suitable almost-periodic KdV data and establish persistence under appropriate perturbations.
\subsection{Research attempt}
\paragraph{Model, date and investigated scope.}
OpenAI GPT-6 Astra (configured model \texttt{gpt-6-astra}), ultra reasoning; generated 11 October 2026. Proved time almost periodicity in uniform spatial norm for absolutely summable Fourier data of the linear Airy equation. Nonlinear KdV and perturbation persistence remain unresolved.

\paragraph{Sources and literature check.}
Read the supplied complete definition, catalogue wording, scope qualifications and source list for this canonical ID. The original-source citations are inherited as supplied references; no fresh primary-source or novelty search was performed for this attempt. No current open-status claim is inferred from this calculation.

\paragraph{A uniform-topology linear benchmark.}Let \(u_0(x)=\sum_{j\ge1}c_je^{i\lambda_jx}\), with real frequencies \(\lambda_j\), \(\sum_j|c_j|<\infty\), and conjugate-paired terms for real data. The linear Airy equation \(u_t+u_{xxx}=0\) has
\[u(x,t)=\sum_jc_je^{i\lambda_jx+i\lambda_j^3t}.
\]
The series converges uniformly in both variables. If \(\sum_j|c_j||\lambda_j|^3<\infty\), termwise differentiation verifies a classical solution. Each finite partial sum is an almost-periodic function of time with values in \(C_b(\mathbb R_x)\), since it is a finite linear combination of fixed vectors with scalar exponential coefficients. Its uniform error is bounded by the coefficient tail independently of time. Hence the full solution is time almost periodic in that Banach norm, using the uniform-approximation definition.
\paragraph{Failure of the immediate nonlinear extension.}For KdV the Fourier coefficient of \(uu_x\) at frequency \(\lambda\) is a convolution containing factors of the differentiated frequency. Absolute summability alone does not control these factors, and repeated nonlinear interactions create new frequencies. Even a local solution with bounded norm would not imply precompactness of its complete time orbit. A proof by approximating general almost-periodic data with periodic data would require estimates uniform for all time, not just on bounded intervals.
\paragraph{Gap and full scope.}This known linear test supplies a topology and a genuine almost-periodicity argument, but no nonlinear global existence or compactness theorem. The first missing step is a globally controlled nonlinear evolution on a suitable almost-periodic function space. The source's small analytic Diophantine results remain supplied historical context, and persistence under suitable nonintegrable perturbations is a separate unresolved part. No universal perturbation class or current-status assertion is inferred.

\paragraph{Verification and adversarial audit.}
Checked Airy phase sign, uniform coefficient tails, classical derivative condition and failure of boundedness or finite-time approximation to imply time almost periodicity. This was an LLM adversarial audit, not independent human verification; no theorem prover was used. Prompts A then B applied to this ID: Checked Airy phase sign, uniform coefficient tails, classical derivative condition and failure of boundedness or finite-time approximation to imply time almost periodicity. Scope remains limited: Construct and control the nonlinear global evolution in a specified almost-periodic topology, uniformly for all time.

Coordinator checked the Airy phase sign, uniform coefficient-tail approximation and the additional derivative-summability requirement. This additional check was model-assisted, not independent human verification.

\paragraph{Tools, process and contributions.}
Python 3 and local shell used to read supplied files and serialize this manuscript; mathematics derived symbolically in the text, with no numerical experiment or formal proof assistant. Contribution prompt A was followed by prompt B. GPT-6 Astra with ultra reasoning drafted the research and performed its stated self-audit; any additional coordinator contribution is identified in the verification and audit. A GPT-6-based Codex coordinator assembled the manuscript and checked repository integration; its exact variant and reasoning configuration were not exposed. No human mathematical contribution, independent human proof check or formal proof certification is supplied. Complete derivations and the consulted source references are included here; private reasoning transcripts are not published.

\paragraph{Final status and first remaining gap.}
\textbf{UNRESOLVED}. Construct and control the nonlinear global evolution in a specified almost-periodic topology, uniformly for all time.

\section{Completion Estimate}
\noindent\textbf{COMPLETION: 0\%}
The calculation is an elementary benchmark or conditional reduction; it establishes no claimed new portion of the complete source target.

\subsection{Sources}
\begin{itemize}
\item[\hypertarget{source:8100001:1}{S1}] ULAM.AI and the UnsolvedMath Contributors, UnsolvedMath, record 8100001 (AMR-080-0001). Catalogue curation under CC BY 4.0; definitions and English exposition added locally.
\url{https://huggingface.co/datasets/ulamai/UnsolvedMath}.
\item[\hypertarget{source:8100001:2}{S2}] Percy Deift, Some Open Problems in Random Matrix Theory and the Theory of Integrable Systems. II, SIGMA 13 (2017), 016, Problem 1.
\url{https://arxiv.org/abs/1703.04931}.
\item[\hypertarget{source:8100001:3}{S3}] Primary full-text formulation and 2017 historical progress discussion.
\url{https://arxiv.org/html/1703.04931v1}.
\end{itemize}
\label{end:8100001}
\end{document}
